% Section: Discussion and Limitations
\section{Discussion}
\label{sec:discussion}

\subsection{Key Findings}
The oracle experiment establishes that hidden courtesy is a
decision-relevant latent variable: knowing whether a merge vehicle is
cooperative or non-cooperative reduces collisions by 8 percentage points
with moderate effect size ($|h|{=}0.230$).
The benefit is concentrated in non-cooperative episodes, where the merge
vehicle narrows the safe-action margin; cooperative vehicles create gaps
large enough that the uniform-prior heuristic already yields safely.

The POMCP result indicates that the oracle is an upper bound only for
the heuristic action table, not for the full policy space.
UCT planning over the reward function identifies encounter-specific action
sequences that the hand-designed V-shaped caution rule cannot represent:
POMCP's gain over the oracle-heuristic baseline ($d{=}0.221$) is comparable
to the oracle-heuristic gain over rule ($d{=}0.233$).
This should be read as evidence that planning improves on the hand-designed
action table, not as a claim that POMCP exceeds an optimal policy with perfect
mode knowledge.

\subsection{The Observation Bottleneck}
The evidence suggests that the belief policy does not outperform rule because
the Bayesian filter converges slowly during the critical early interaction
window.
Non-cooperative IDM vehicles briefly decelerate in response to ego
detection before reasserting their high target speed, producing
acceleration signals that temporarily resemble cooperative behaviour;
sustained non-cooperative speed requires roughly 15--20 interaction-window
steps before \emph{mvs} and \emph{mva} accumulate discriminative evidence.
The belief-gain sweep directly rules out action-table expressiveness as
the bottleneck: if the heuristic were too insensitive to the belief signal,
a $10\times$ gain amplification would improve belief-policy performance.
It does not (Table~\ref{tab:gain_sweep}; Figure~\ref{fig:gain_sweep}).
Future designs should target the early-window transient directly---through
recurrent observation encoders, longer pre-interaction observation windows,
or planning architectures that are robust to early-phase ambiguity.

\subsection{Limitations}
All results are from the highway-env \texttt{merge-v0} simulator; no
real-world deployment or human-driver behavioural validation is performed, and
none is claimed---results characterise the simulator only.
Courtesy types are operationalised via IDM/MOBIL parameter ranges and do
not model empirical human psychology or the continuous spectrum of real
driver behaviour.
The $n{=}300$ evaluation is underpowered for the belief-vs.-rule
comparison ($n{\approx}1{,}960$ needed for 80\% power); the null result
should be treated as suggestive.
The 70/15/15 dataset split is assigned by episode index at the same global
seed; sequentially generated episodes may share correlated traffic densities,
reducing the effective sample size for any model trained on the dataset.
An empirical lag-1 autocorrelation check on episode-level covariates shows
that traffic density and observation noise have lag-1 correlations $<0.03$,
indicating that sequential correlation is negligible in practice; the hidden
courtesy type shows lag-1 correlation $-0.50$, reflecting the generator's
near-perfect alternation that ensures balanced courtesy classes and should be
documented because it affects split construction.
The dataset was generated with \texttt{belief\_policy} as the ego,
under-representing high-risk states relative to a random ego; the seed-17
matched policy evaluation is unaffected, but researchers training
learning-based planners on the released dataset should account for this
trajectory-distribution bias (see \S\ref{sec:benchmark}).
The POMCP generative model uses mean IDM parameters per mode rather than
sampling from the full injected distributions, introducing model error
whose sensitivity is not studied here.
A rollout-isolation run replaces the conservative POMCP rollout trigger
(SLOWER when $d{<}15$\,m or TTC\,$<$\,6\,s) with the heuristic emergency
trigger (TTC\,$<$\,3\,s or gap\,$<$\,8\,m) and obtains the same collision
rate with nearly identical reward (Table~\ref{tab:pomcp_rollout_ablation}).
This reduces, but does not eliminate, the rollout-confound concern because
the POMCP generative model and planning reward still differ from the
environment reward.
The planning reward uses a narrower collision zone ($d{<}1.5$\,m) and
omits the environment's survival bonus, creating a planning--evaluation
reward mismatch that may bias action selection during episodes in which the
ego is close to but not inside the $1.5$\,m threshold.
The observation model's product-of-Gaussians independence assumption is
imperfectly satisfied (largest residual correlation: \emph{mvs}/\emph{mva},
$r{=}0.399$) and the 5.2:1 calibration step imbalance increases
cooperative parameter uncertainty.
A held-out full-covariance Gaussian improves offline belief metrics
(Table~\ref{tab:covariance_ablation}); closed-loop reruns with this likelihood
remain an important benchmark extension.
A naive majority-class (non-cooperative) classifier achieves 84\% step-level
accuracy on calibration window steps; the three-feature model's 98.3\%
confirms that discriminative information is drawn from the features rather
than the class prior.
The 12-second episode cap ($60{\times}0.2$\,s) is sufficient for the
majority of merge interactions at the benchmark traffic density, but late-starting
merge events where the merge vehicle enters the interaction window
after step 40 experience fewer than 4\,s of available observation, reducing
belief accuracy for those episodes.
The QMDP approximation (16-state discretisation, $n{=}50$ training
episodes) achieved 34--35\% collision rate with 50\% Q-table state
coverage and is excluded from the primary evaluation.
